\documentclass[12pt,reqno]{article}

\usepackage[usenames]{color}
\usepackage{amssymb}
\usepackage{amsmath}
\usepackage{amsthm}
\usepackage{amsfonts}
\usepackage{amscd}
\usepackage{graphicx}
\usepackage{mathtools}

\mathtoolsset{showonlyrefs}


%\usepackage{showkeys}
\usepackage[colorlinks=true,
linkcolor=webgreen,
filecolor=webbrown,
citecolor=webgreen]{hyperref}

\definecolor{webgreen}{rgb}{0,.5,0}
\definecolor{webbrown}{rgb}{.6,0,0}

\usepackage{color}
\usepackage{fullpage}
\usepackage{float}

%\usepackage{psfig}
\usepackage{graphics}
\usepackage{latexsym}
%\usepackage{epsf}
%\usepackage{breakurl}
%\usepackage{accents}

\setlength{\textwidth}{6.5in}
\setlength{\oddsidemargin}{.1in}
\setlength{\evensidemargin}{.1in}
\setlength{\topmargin}{-.1in}
\setlength{\textheight}{8.4in}

\newcommand{\seqnum}[1]{\href{https://oeis.org/#1}{\rm \underline{#1}}}

\DeclareMathOperator{\arctanh}{arctanh}
\DeclareMathOperator{\Li}{Li}
\DeclareMathOperator{\Cl}{Cl}

\newcommand{\blue}[1]{{\color{blue}#1}}
\newcommand{\red}[1]{{\color{red}#1}}
\newcommand{\magenta}[1]{{\color{magenta}#1}}

\newcommand{\bm}[1]{{\boldmath{\text{$#1$}}}}

%{\boldmath{\text{$ab$}}}


\begin{document}

\theoremstyle{plain}
\newtheorem{theorem}{Theorem}
\newtheorem{corollary}[theorem]{Corollary}
\newtheorem{proposition}{Proposition}
\newtheorem{lemma}{Lemma}
%\newtheorem*{example}{Examples}
\newtheorem{example}{Example}
\newtheorem{remark}{Remark}

% \numberwithin{equation}{section}
% \numberwithin{example}{section}
% \numberwithin{theorem}{section}
% \numberwithin{remark}{section}
% \numberwithin{proposition}{section}
% \numberwithin{lemma}{section}



\newcommand{\braces}{\genfrac{\lbrace}{\rbrace}{0pt}{}}

\begin{center}

\vskip 1cm{\large\bf  
New Generalizations of Two Ramanujan Series for $1/\pi$}
\vskip 1cm
\large
Kunle Adegoke \\ 
Department of Physics and Engineering Physics\\
Obafemi Awolowo University\\
220005 Ile-Ife \\
Nigeria \\
\href{mailto:adegoke00@gmail.com }{\tt adegoke00@gmail.com}\\

\end{center}

\vskip .2 in

\begin{abstract}
By utilizing two hypergeometric summation identities, we extend two \mbox{well-known} Ramanujan series for $1/\pi$ by making each of them a member of an infinite family of series. We also find the corresponding families involving harmonic numbers and odd harmonic numbers. 
\end{abstract}

\noindent 2020 {\it Mathematics Subject Classification}: Primary 33C20; Secondary 11Y60. 

\noindent \emph{Keywords:} Ramanujan series, Ramanujan-like series, harmonic number, odd harmonic number, hypergeometric series.

\section{Introduction}
Let $\left( x \right)_n  = \prod_{k = 0}^{n - 1} {\left( {x + k} \right)} $, for $x$ a complex number and $n$ a non-negative integer. Our main purpose in this note is to derive the following families of Ramanujan-like series for $1/\pi$:
\begin{align}\label{gen1}
&\sum_{k = 0}^\infty  {\frac{{( - 1)^k }}{{k!^3 }}\frac{{\left( {\frac{1}{2}} \right)_k^3 \left( {4k + 1 - 2\ell} \right)}}{{\prod_{j = 1}^\ell {\left( {2k - 2j + 1} \right)} \prod_{j = 1}^m {\left( {2k - 2j + 1} \right)} \prod_{j = 1}^n {\left( {2k - 2j + 1} \right)} \prod_{j = 1}^{m - \ell} {\left( {k + j} \right)} \prod_{j = 1}^{n - \ell} {\left( {k + j} \right)} }}}\nonumber\\
&\qquad  = \frac{{( - 1)^{m + n} }}{{\left( {\frac{1}{2}} \right)_m \left( {\frac{1}{2}} \right)_n \left( {\frac{1}{2}} \right)_{m + n - \ell} 2^{\ell + m + n - 1} }}\,\frac{1}{\pi }
\end{align}
and
\begin{equation}\label{gen2}
\sum_{k = 0}^\infty  {\frac{{\left( {\frac{1}{2}} \right)_k^3 }}{{4^k k!^3 }}\frac{{\prod_{j = 1}^{2n} {\left( {2k + j} \right)} }}{{\prod_{j = 1}^n {\left( {k + j} \right)^2 \left( {2k - 2j + 1} \right)} }}\left( {6k + 1 + 2n} \right)}  = \frac{{( - 1)^n\,2^{n+2} }}{{\left( {\frac{1}{2}} \right)_n }}\, \frac1{\pi };
\end{equation}
valid for $\ell$, $m$, and $n$ non-negative integers.

Identities~\eqref{gen1} and~\eqref{gen2} are respective generalizations of the Ramanujan series~\cite{ramanujan14}:
\begin{equation}\label{bauer}
\sum_{k = 0}^\infty  {( - 1)^k \frac{{\left( \frac12 \right)_k^3 }}{{k!^3 }}\left( {4k + 1} \right)}  = \frac{2}{\pi }
\end{equation}
and
\begin{equation}
\sum_{k = 0}^\infty  {\frac{{\left( {\frac{1}{2}} \right)_k^3 }}{{4^k k!^3 }}\left( {6k + 1} \right)}  = \frac{4}{\pi }.
\end{equation}

When $\ell=0=m$, identity~\eqref{gen1} reduces to
\begin{equation*}
\sum_{k = 0}^\infty  {( - 1)^k \frac{{\left( {\frac{1}{2}} \right)_k^3 }}{{k!^3 }}\frac{{\left( {4k + 1} \right)}}{{\prod_{j = 1}^n {\left( {2k - 2j + 1} \right)\left( {k + j} \right)} }}}  = \frac{{( - 1)^n }}{{\left( {\frac{1}{2}} \right)_n^2 2^{n - 1} }}\,\frac{1}{\pi },
\end{equation*}
which was derived by Levrie~\cite[Theorem 7]{levrie10} using Fourier-Legendre theory. In fact we will see that identities~\eqref{gen1} and~\eqref{gen2} are particular cases of more general results.

Let $H_j$ be the $j^{\rm th}$ harmonic number and let $O_j$ be the $j^{\rm th}$ odd harmonic number. We will also derive the following series involving harmonic and odd harmonic numbers:
\begin{equation}\label{di4jyj5}
\sum_{k = 0}^\infty  {\frac{{( - 1)^k \left( {\frac{1}{2}} \right)_k^3 \left( {4k + 1} \right)\left( {2O_{k - n}  + H_{k + n} } \right)}}{{k!^3 \prod_{j = 1}^n {\left( {k + j} \right)\left( {2k - 2j + 1} \right)} }}}  = \frac{{( - 1)^n }}{{2^{n - 2} }}\frac{1}{{\left( {\frac{1}{2}} \right)_n^2 }}\left( {2O_n  - \ln 2} \right)\,\frac{1}{\pi }
\end{equation}
and
\begin{align}\label{z75f5e1}
&\sum_{k = 0}^\infty  {\frac{{\left( {\frac{1}{2}} \right)_k^3 }}{{4^k k!^3 }}\frac{{\prod_{j = 1}^{2n} {\left( {2k + j} \right)} }}{{\prod_{j = 1}^n {\left( {k + j} \right)^2 \left( {2k - 2j + 1} \right)} }}\left( {\left( {6k + 1 + 2n} \right)\left( {2O_{k - n}  - 2O_{k + n}  + H_{k + n} } \right) - 2} \right)}\nonumber\\
&\qquad  = ( - 1)^n\,\frac{{2^{n + 3} }}{{\left( {\frac{1}{2}} \right)_n }}\left( {O_n  - \ln 2} \right)\frac{1}{\pi }.
\end{align}

Identity~\eqref{di4jyj5} at $n=0$ reduces to the harmonic number complement of~\eqref{bauer}:
\begin{equation}
\sum_{k = 0}^\infty  {( - 1)^k \frac{{\left( {\frac{1}{2}} \right)_k^3 }}{{k!^3 }}\left( {4k + 1} \right)H_{2k} }  =  - \frac{2\ln 2}{\pi },
\end{equation}
derived by Hou et al.~\cite{hou24} and also Campbell~\cite{campbell26}.

The special case $n=0$ of identity~\eqref{z75f5e1}, namely,
\begin{equation}
\sum_{k = 0}^\infty  {\frac{{\left( {\frac{1}{2}} \right)_k^3 }}{{4^k k!^3 }}\left( {\left( {6k + 1} \right)H_k  - 2} \right)}  =  - \frac{{8\ln 2}}{\pi },
\end{equation}
was derived by Hou et al.~\cite{hou24}.

The now famous identity~\eqref{bauer} was discovered in 1859 by Bauer~\cite{bauer} via a Fourier-Legendre expansion, and rediscovered in 1914 by Ramanujan who included it in his famous first letter to G.~H.~Hardy. Early proofs relied on classical techniques like swapping limit operations while later approaches utilized advanced tools like Dougall's bilateral hypergeometric theorem. For a survey of the \mbox{Bauer-Ramanujan} series~\eqref{bauer} and related results and analyses, the reader is referred to the recent article by Campbell and Levrie~\cite{campbell24}.

Let $\Gamma(z)$ be Euler's Gamma function. Generalized binomial coefficients are defined for complex numbers $r$ and $s$ by
\begin{equation}\label{y89d722}
\binom rs= \frac{{\Gamma (r + 1)}}{{\Gamma (s + 1)\Gamma (r - s + 1)}}.
\end{equation}
The {\it extended} Pochhammer relation $(x)_y$ is defined for complex numbers $x$ and $y$ that are not non-positive integers and such that $x+y\ne0$ by
\begin{equation}\label{ex_pochhammer}
\left( x \right)_y  = \frac{{\Gamma \left( {x + y} \right)}}{{\Gamma \left( x \right)}} = \binom{{x + y - 1}}{y}\,\Gamma \left( {y + 1} \right).
\end{equation}
If, in particular,  $y=n$ is a non-negative integer, then we have the Pochhammer symbol
\begin{equation}
\left( x \right)_n  = \prod_{k = 0}^{n - 1} {\left(x + k\right)}  = \binom{{x + n - 1}}{n}\,n!
\end{equation}

Harmonic numbers, $H_j$, and odd harmonic numbers, $O_j$, are defined for non-negative integers $j$ by
\begin{equation}
H_j=\sum_{k=1}^j\frac1k,\quad O_j=\sum_{k=1}^j\frac1{2k-1},\quad H_0=0=O_0.
\end{equation}
The identity
\begin{equation}\label{nerd5x1}
H_{n-1/2}=2O_n-2\ln 2,
\end{equation}
extends harmonic numbers to half-integer arguments and is a consequence of the link between harmonic numbers, odd harmonic numbers and the digamma function.

\section{Required results}
\begin{lemma}
If $k$ and $x$ are complex numbers, then
\begin{equation}\label{yahnm05}
\binom{{k}}{{x + 1/2}} = \frac{2}{\pi }\,\frac{1}{{\left( {2x + 1} \right)}}\frac{{x!k!}}{{\left( {\frac{1}{2}} \right)_x \left( {\frac{1}{2}} \right)_k }}\binom{{k - 1/2}}{x}.
\end{equation}
\end{lemma}
\begin{proof}
The identity
\begin{equation}\label{w7han3w}
\binom{{a - b}}{c} = \binom{{a - c}}{b}\binom{{a}}{c}\binom{{a}}{b}^{ - 1} ,
\end{equation}
with $a=k$, $b=1/2$ and $c=k-1/2-x$ gives
\begin{equation}
\binom{{k - 1/2}}{x} = \binom{{x + 1/2}}{{1/2}}\binom{{k}}{{x + 1/2}}\binom{{k}}{{1/2}}^{ - 1}; 
\end{equation}
and hence~\eqref{yahnm05} since
\begin{equation}
\binom{{x + 1/2}}{{1/2}} = \frac{{2x + 1}}{{x!}}\left( {\frac{1}{2}} \right)_x 
\end{equation}
and
\begin{equation}
\binom{{k}}{{1/2}} = \frac{2}{\pi }\,\frac{{k!}}{{\left( {\frac{1}{2}} \right)_k }}.
\end{equation}
\end{proof}
\begin{lemma}\label{w00yljy}
If $x$ is a complex number that is not a negative integer, then
\begin{equation}\label{l5fmdri}
\left( { - x + \frac{1}{2}} \right)_k  = \frac{{\left( {\frac{1}{2}} \right)_x \left( {\frac{1}{2}} \right)_k }}{{x!}}\,\binom{{k - 1/2}}{x}^{ - 1} \cos \left( \pi x \right)
\end{equation}
and
\begin{equation}\label{bjv8z0f}
\left( {x + \frac{1}{2}} \right)_k  = \left( {\frac{1}{2}} \right)_k \binom{{2k + 2x}}{{2x}}\binom{{k + x}}{x}^{ - 1}.
\end{equation}
In particular, if $n$ is a non-negative integer, then
\begin{equation}\label{mq7ny8m}
\left( { - n + \frac{1}{2}} \right)_k  = ( - 1)^n \frac{{2^n \left( {\frac{1}{2}} \right)_n \left( {\frac{1}{2}} \right)_k }}{{\prod_{j = 1}^n {\left( {2k - 2j + 1} \right)} }}
\end{equation}
and
\begin{equation}\label{t0khqm8}
\left( {n + \frac{1}{2}} \right)_k  = \frac{1}{{2^{2n} }}\frac{{\left( {\frac{1}{2}} \right)_k }}{{\left( {\frac{1}{2}} \right)_n }}\frac{{\prod_{j = 1}^{2n} {\left( {2k + j} \right)} }}{{\prod_{j = 1}^n {\left( {k + j} \right)} }}.
\end{equation}
\end{lemma}
\begin{proof}
We prove~\eqref{l5fmdri}.
\begin{equation}
\left( { - x + \frac{1}{2}} \right)_k  = k!\,\binom{{k - x - 1/2}}{k}.
\end{equation}
Now, using~\eqref{w7han3w} with $a=k$, $b=x+1/2$ and $c=k$, write
\begin{equation}
\binom{{k - x - 1/2}}{k} = \binom{{0}}{{x + 1/2}}\binom{{k}}{k}\binom{{k}}{{x + 1/2}}^{ - 1}  = \frac{2}{\pi }\,\frac{{\cos \left( \pi x\right)}}{{2x + 1}}\binom{{k}}{{x + 1/2}}^{ - 1},
\end{equation}
since
\begin{equation}
\binom{{0}}{s}
=\begin{cases}
 \dfrac{{\sin \left( {\pi s} \right)}}{{\pi s}},&\text{if $s\ne0$;} \\ 
 1,&\text{if $s=0$.} \\ 
 \end{cases}
\end{equation}
Identity~\eqref{l5fmdri} now follows by the use of~\eqref{yahnm05}.
\end{proof}
\begin{lemma}
We have
\begin{equation}
\frac{d}{{dx}}\left( x \right)_y  = \left( x \right)_y \left( {H_{x + y - 1}  - H_{x - 1} } \right).
\end{equation}
\end{lemma}

\section{Generalizations of Ramanujan series for $1/\pi$}
\begin{lemma}\label{rutep3z}
If $k$ is a non-negative integer, then
\begin{equation}
\lim_{d\to\infty}{\frac{{\left( d \right)_k }}{{\left( { - d} \right)_k }}} = ( - 1)^k .
\end{equation}
\end{lemma}
\begin{proof}
Now
\begin{equation}
\left( { - d} \right)_k  = k!\binom{{ - d + k - 1}}{k} = ( - 1)^k k!\binom{{d}}{k}.
\end{equation}
But
\begin{equation}
\left( d \right)_k  = k!\binom{{d + k - 1}}{k} \sim k!\binom{{d}}{k}\text{ as $d\to\infty$}.
\end{equation}
Hence the result.
\end{proof}
\begin{lemma}\label{z2u9bwg}
If $a$ and $k$ are complex numbers, then
\begin{equation}
\frac{{\left( {1 + \frac{a}{2}} \right)_k }}{{\left( {\frac{a}{2}} \right)_k }} = \frac{{2k + a}}{a}.
\end{equation}
\end{lemma}
\begin{lemma}
For suitably bounded complex numbers $a$, $b$, and $c$, the following identity holds:
\begin{equation}\label{rbz4noe}
\sum_{k = 0}^\infty  {\frac{{( - 1)^k }}{{k!}}\frac{{\left( a \right)_k \left( b \right)_k \left( c \right)_k }}{{\left( {1 + a - b} \right)_k \left( {1 + a - c} \right)_k }}\left( {2k + a} \right)}  = \frac{{a\Gamma \left( {1 + a - b} \right)\Gamma \left( {1 + a - c} \right)}}{{\Gamma \left( {1 + a} \right)\Gamma \left( {1 + a - b - c} \right)}}.
\end{equation}
In particular,
\begin{equation}\label{q0xvexn}
\sum_{k = 0}^\infty  {\frac{{( - 1)^k }}{{k!^2 }}\frac{{\left( a \right)_k^2 \left( c \right)_k }}{{\left( {1 + a - c} \right)_k }}\left( {2k + a} \right)}  = \frac{{a\Gamma \left( {1 + a - c} \right)}}{{\Gamma \left( {1 + a} \right)\Gamma \left( {1 - c} \right)}}
\end{equation}
and
\begin{equation}\label{nw1xvnr}
\sum_{k = 0}^\infty  {( - 1)^k \frac{{\left( a \right)_k^3 }}{{k!^3 }}\left( {2k + a} \right)}  = \frac{{\sin \left( {\pi a} \right)}}{\pi }.
\end{equation}
\end{lemma}
\begin{proof}
Consider Dougall's well-poised hypergeometric series~\cite[p.~27]{bailey35}:
\begin{align}\label{ym6fvmk}
&\sum_{k = 0}^\infty  {\frac{1}{{k!}}\frac{{\left( a \right)_k \left( {1 + a/2} \right)_k \left( b \right)_k \left( c \right)_k \left( d \right)_k }}{{\left( {a/2} \right)_k \left( {1 + a - b} \right)_k \left( {1 + a - c} \right)_k \left( {1 + a - d} \right)_k }}}\nonumber\\
&\qquad  = \frac{{\Gamma \left( {1 + a - b} \right)\Gamma \left( {1 + a - c} \right)\Gamma \left( {1 + a - d} \right)\Gamma \left( {1 + a - b - c - d} \right)}}{{\Gamma \left( {1 + a} \right)\Gamma \left( {1 + a - b - c} \right)\Gamma \left( {1 + a - b - d} \right)\Gamma \left( {1 + a - c - d} \right)}},
\end{align}
which holds for suitably bounded complex numbers $a$, $b$, $c$ and $d$ or $\Re(1+a-b-c-d)>0$ if each of the four numbers has a finite magnitude.

Letting $d$ approach infinity in~\eqref{ym6fvmk} gives~\eqref{rbz4noe} by application of Lemmata~\ref{rutep3z} and~\ref{z2u9bwg}.
\end{proof}

\begin{remark}
Campbell~\cite{campbell26} also derived~\eqref{q0xvexn} using Mishev's transform. Identity~\eqref{nw1xvnr} is Dougall's identity~\cite[Equation (16)]{dougall06}.
\end{remark}
The identity stated in Lemma~\ref{edacukd} corresponds to setting $a=1/2=b$ in~\eqref{rbz4noe} and was derived by Hou et al.~\cite{hou24}. 

\begin{lemma}\label{edacukd}
If $c$ is a complex number, then
\begin{equation}\label{o68u3j7}
\sum_{k = 0}^\infty  {( - 1)^k \frac{{\left( {\frac{1}{2}} \right)_k^2 \left( c \right)_k }}{{k!^2 \left( {3/2 - c} \right)_k }}\left( {4k + 1} \right)}  = \frac{2}{{\sqrt \pi  }}\,\frac{{\Gamma \left( {3/2 - c} \right)}}{{\Gamma \left( {1 - c} \right)}}= \frac{2}{{\sqrt \pi  }}(1-c)_{1/2}.
\end{equation}
\end{lemma}

\begin{remark}
The special case of~\eqref{o68u3j7} with $c$ a non-negative integer was proved by Ekhad and Zeilberger~\cite{ekhad94} using the $WZ$ technique.
\end{remark}

\begin{theorem}
If $p$, $q$, and $r$ are complex numbers that are not half-integers and such that $q-p$ and $r-p$ are not non-positive integer, then
\begin{align}
&\sum_{k = 0}^\infty  {( - 1)^k \frac{{\left( {\frac{1}{2}} \right)_k^3 }}{{k!^3 }}\,\frac{{\left( {4k + 1 - 2p} \right)}}{{\binom{{k - 1/2}}{p}\binom{{k - 1/2}}{q}\binom{{k - 1/2}}{r}\binom{{k + q - p}}{k}\binom{{k + r - p}}{k}}}}\\
&\qquad= 2\,\frac{{\left( {\frac{1}{2} - p} \right)_{q + \frac{1}{2}} }}{{\left( {r - p} \right)_{q + \frac{1}{2}} }}\,\frac{{p!q!r!}}{{\left( {\frac{1}{2}} \right)_p \left( {\frac{1}{2}} \right)_q \left( {\frac{1}{2}} \right)_r }}\,\frac{{r - p}}{{\cos \left( {\pi p} \right)\cos \left( {\pi q} \right)\cos \left( {\pi r} \right)}} 
\end{align}
In particular,
\begin{equation}
\sum_{k = 0}^\infty  {( - 1)^k \frac{{\left( {\frac{1}{2}} \right)_k^3 }}{{k!^3 }}\frac{{\left( {4k + 1} \right)}}{{\binom{{k + r}}{k}\binom{{k - 1/2}}{r}}}}  = \frac{2}{\pi }\,\frac{{r!^2 }}{{\left( {\frac{1}{2}} \right)_r^2 }}\,\frac{1}{{\cos \left( \pi r \right)}}.
\end{equation}
\end{theorem}
\begin{proof}
Set $a=1/2-p$, $b=1/2-q$ and $c=1/2-r$ in~\eqref{rbz4noe} to obtain
\begin{align}\label{h7933kq}
&\sum_{k = 0}^\infty  {\frac{{( - 1)^k }}{{k!}}\frac{{\left( {1/2 - p} \right)_k \left( {1/2 - q} \right)_k \left( {1/2 - r} \right)_k }}{{\left( {1 - p + q} \right)_k \left( {1 - p + r} \right)_k }}\left( {4k + 1 - p} \right)}\nonumber\\
&\qquad  = \frac{{2\,\Gamma \left( {1 - p + q} \right)\Gamma \left( {1 - p + r} \right)}}{{\Gamma \left( {1/2 - p} \right)\Gamma \left( {1/2 - p + q + r} \right)}} .
\end{align}
Use~\eqref{l5fmdri} and note from~\eqref{ex_pochhammer} that
\begin{equation}\label{nly5agg}
\left( {1 + x} \right)_k  = k!\binom{{x + k}}{k},
\end{equation}
\begin{equation}
\frac{{\Gamma \left( {1 - p + q} \right)}}{{\Gamma \left( {\frac{1}{2} - p} \right)}} = \left( {\frac{1}{2} - p} \right)_{\frac{1}{2} + q} ,
\end{equation}
and
\begin{equation}
\frac{{\Gamma \left( {1 - p + r} \right)}}{{\Gamma \left( {\frac{1}{2} - p + q + r} \right)}} = \frac{{r - p}}{{\left( {r - p} \right)_{\frac{1}{2} + q} }}.
\end{equation}
\end{proof}
\begin{corollary}
If $\ell$, $m$ and $n$ are non-negative integers, then
\begin{align*}
&\sum_{k = 0}^\infty  {\frac{{( - 1)^k }}{{k!^3 }}\frac{{\left( {\frac{1}{2}} \right)_k^3 \left( {4k + 1 - 2\ell} \right)}}{{\prod_{j = 1}^\ell {\left( {2k - 2j + 1} \right)} \prod_{j = 1}^m {\left( {2k - 2j + 1} \right)} \prod_{j = 1}^n {\left( {2k - 2j + 1} \right)} \prod_{j = 1}^{m - \ell} {\left( {k + j} \right)} \prod_{j = 1}^{n - \ell} {\left( {k + j} \right)} }}}\\
&\qquad  = \frac{{( - 1)^{m + n} }}{{\left( {\frac{1}{2}} \right)_m \left( {\frac{1}{2}} \right)_n \left( {\frac{1}{2}} \right)_{m + n - \ell} 2^{\ell + m + n - 1} }}\,\frac{1}{\pi }.
\end{align*}
In particular,
\begin{align}
\sum_{k = 0}^\infty  {( - 1)^k \frac{{\left( {\frac{1}{2}} \right)_k^3 }}{{k!^3 }}\frac{{\left( {4k + 1} \right)}}{{\prod_{j = 1}^n {\left( {2k - 2j + 1} \right)\left( {k + j} \right)} }}}  &= \frac{{( - 1)^n }}{{\left( {\frac{1}{2}} \right)_n^2 2^{n - 1} }}\,\frac{1}{\pi },\\
\sum_{k = 0}^\infty  {( - 1)^k \frac{{\left( {\frac{1}{2}} \right)_k^3 }}{{k!^3 }}\frac{{\left( {4k + 1 - 2n} \right)}}{{\prod_{j = 1}^n {\left( {2k - 2j + 1} \right)^3 } }}}  &= \frac{1}{{\left( {\frac{1}{2}} \right)_n^3 2^{3n - 1} }}\,\frac{1}{\pi },
\end{align}
\begin{align}
&\sum_{k = 0}^\infty  {( - 1)^k \frac{{\left( {\frac{1}{2}} \right)_k^3 }}{{k!^3 }}\frac{{\left( {4k + 1 - 2m} \right)}}{{\prod_{j = 1}^m {\left( {2k - 2j + 1} \right)^2 } \prod_{j = 1}^n {\left( {2k - 2j + 1} \right)} \prod_{j = 1}^{n - m} {\left( {k + j} \right)} }}}\nonumber\\
&\qquad  = \frac{{( - 1)^{m + n} }}{{\left( {\frac{1}{2}} \right)_m \left( {\frac{1}{2}} \right)_n^2 2^{2m + n - 1} }}\,\frac{1}{\pi },
\end{align}
and
\begin{align}
&\sum_{k = 0}^\infty  {( - 1)^k \frac{{\left( {\frac{1}{2}} \right)_k^3 }}{{k!^3 }}\frac{{\left( {4k + 1 - 2\ell} \right)}}{{\prod_{j = 1}^\ell {\left( {2k - 2j + 1} \right)} \prod_{j = 1}^n {\left( {2k - 2j + 1} \right)^2 } \prod_{j = 1}^{n - \ell} {\left( {k + j} \right)^2 } }}}\nonumber\\
&\qquad  = \frac{1}{{\left( {\frac{1}{2}} \right)_{2n - l} \left( {\frac{1}{2}} \right)_n^2 2^{2n + \ell - 1} }}\,\frac{1}{\pi }.
\end{align}
\end{corollary}
\begin{proof}
Set $p=\ell$, $q=m$ and $r=n$ in~\eqref{h7933kq} and use~\eqref{mq7ny8m} and~\eqref{nly5agg}. Note that
\begin{equation}
n!\binom{{k + n}}{n} = \prod_{j = 1}^n {\left( {k + j} \right)},\text{ etc} ,
\end{equation}
\begin{equation}
\Gamma \left( {\frac{1}{2} - u} \right) = \frac{{( - 1)^u }}{{\left( {\frac{1}{2}} \right)_u }}\,\frac{{\sqrt \pi  }}{{u!}},
\end{equation}
and
\begin{equation}
\Gamma \left( {\frac{1}{2} + u} \right) = \left( {\frac{1}{2}} \right)_u \sqrt \pi .
\end{equation}
\end{proof}
\begin{theorem}
If $r$ is a complex number that is not a half-integer, then
\begin{equation}\label{r1kycoj}
\sum_{k = 0}^\infty  {( - 1)^k \frac{{\left( {\frac{1}{2}} \right)_k^3 }}{{k!^3 }}\frac{{\left( {4k + 1} \right)\left( {2O_{k - r}  + H_{k + r} } \right)}}{{\binom{{k + r}}{k}\binom{{k - 1/2}}{r}}}}  = \frac{2}{\pi }\,\frac{{r!^2 }}{{\left( {\frac{1}{2}} \right)_r^2 }}\,\frac{{\left( {4O_r  - 2\ln 2} \right)}}{{\cos \left( \pi r \right)}}.
\end{equation}
\end{theorem}
\begin{proof}
Differentiate~\eqref{o68u3j7} with respect to $c$ and replace $c$ with $1/2-r$ to get
\begin{equation}\label{hr8o23s}
\sum_{k = 0}^\infty  {\frac{{( - 1)^k \left( {\frac{1}{2}} \right)_k^2 \left( {\frac{1}{2} - r} \right)_k \left( {4k + 1} \right)\left( {H_{k - r - 1/2}  + H_{k + r} } \right)}}{{k!^2 \left( {r + 1} \right)_k }}}  = \frac{2}{{\sqrt \pi  }}\left( {r + \frac{1}{2}} \right)_{1/2} \left( {H_{ - r - 1/2}  + H_{r - 1/2} } \right),
\end{equation}
from which~\eqref{r1kycoj} follows after using~\eqref{nerd5x1},~\eqref{l5fmdri},~\eqref{nly5agg} and the fact that
\begin{equation}\label{g20r5hg}
\left( {r + \frac{1}{2}} \right)_{\frac{1}{2}}  = \frac{{\Gamma \left( {r + 1} \right)}}{{\Gamma \left( {r + \frac{1}{2}} \right)}} = \frac{{r!}}{{\left( {\frac{1}{2}} \right)_r \sqrt \pi  }}.
\end{equation}
\end{proof}
\begin{corollary}
If $n$ is a non-negative integer, then
\begin{equation*}
\sum_{k = 0}^\infty  {\frac{{( - 1)^k \left( {\frac{1}{2}} \right)_k^3 \left( {4k + 1} \right)\left( {2O_{k - n}  + H_{k + n} } \right)}}{{k!^3 \prod_{j = 1}^n {\left( {k + j} \right)\left( {2k - 2j + 1} \right)} }}}  = \frac{{( - 1)^n }}{{2^{n - 2} }}\frac{1}{{\left( {\frac{1}{2}} \right)_n^2 }}\left( {2O_n  - \ln 2} \right)\,\frac{1}{\pi }.
\end{equation*}
In particular,
\begin{equation*}
\sum_{k = 0}^\infty  {( - 1)^k \frac{{\left( {\frac{1}{2}} \right)_k^3 }}{{k!^3 }}\left( {4k + 1} \right)H_{2k} }  =  - \frac{2}{\pi }\ln 2.
\end{equation*}
\end{corollary}
\begin{proof}
Set $r=n$ a non-negative integer in~\eqref{hr8o23s}.
\end{proof}

Lemma~\ref{nnoqyfq} was derived by Hou et al.~\cite{hou24} using a hypergeometric transformation formula due to Chu and Zhang~\cite[Theorem 9]{chu14}; and facilitates the derivation of~\eqref{gen2}.

\begin{lemma}\label{nnoqyfq}
If $c$ is a complex number, then
\begin{equation}\label{k1x81g9}
\sum_{k = 0}^\infty  {\frac{1}{{4^k }}\frac{{\left( {\frac{1}{2}} \right)_k \left( c \right)_k \left( {1 - c} \right)_k }}{{k!^2 \left( {3/2 - c} \right)_k }}\,\frac{{3k - c + 1}}{2}}  = \frac{1}{{\sqrt \pi  }}\,(1-c)_{1/2}.
\end{equation}
\end{lemma}
\begin{theorem}
If $r$ is a complex number that is not a half-integer, then
\begin{equation}
\sum_{k = 0}^\infty  {\frac{{\left( {\frac{1}{2}} \right)_k^3 }}{{4^k k!^3 }}\frac{{\binom{{2k + 2r}}{{2r}}}}{{\binom{{k + r}}{r}^2 \binom{{k - 1/2}}{r}}}\left( {6k + 1 + 2r} \right)}  = \frac{4}{\pi }\frac{{r!^2 }}{{\left( {\frac{1}{2}} \right)_r^2 }}\,\frac{{1 }}{{\cos\left( {\pi r} \right)}}.
\end{equation}
\end{theorem}
\begin{proof}
Set $c=1/2-r$ in~\eqref{k1x81g9} to obtain
\begin{equation}\label{w4w5kv4}
\sum_{k = 0}^\infty  {\frac{1}{{4^k }}\frac{{\left( {\frac{1}{2}} \right)_k \left( {\frac{1}{2} - r} \right)_k \left( {\frac{1}{2} + r} \right)_k }}{{k!^2 \left( {r + 1} \right)_k }}\left( {6k + 1 + 2r} \right)}  = \frac{4}{{\sqrt \pi  }}\left( {r + \frac{1}{2}} \right)_{1/2} .
\end{equation}
Now use~\eqref{l5fmdri},~\eqref{bjv8z0f},~\eqref{nly5agg} and~\eqref{g20r5hg}.
\end{proof}
\begin{corollary}
If $n$ is a non-negative integer, then
\begin{equation*}
\sum_{k = 0}^\infty  {\frac{{\left( {\frac{1}{2}} \right)_k^3 }}{{4^k k!^3 }}\frac{{\prod_{j = 1}^{2n} {\left( {2k + j} \right)} }}{{\prod_{j = 1}^n {\left( {k + j} \right)^2 \left( {2k - 2j + 1} \right)} }}\left( {6k + 1 + 2n} \right)}  = \frac{{( - 1)^n\,2^{n+2} }}{{\left( {\frac{1}{2}} \right)_n }}\, \frac1{\pi }.
\end{equation*}
In particular,
\begin{equation}
\sum_{k = 0}^\infty  {\frac{{\left( {\frac{1}{2}} \right)_k^3 }}{{4^k k!^3 }}\frac{{\left( {2k + 1} \right)^2 }}{{\left( {k + 1} \right)\left( {2k - 1} \right)}}}  =  - \frac{8}{{3\pi }}.
\end{equation}
\end{corollary}
\begin{proof}
Set $r=n$ in~\eqref{w4w5kv4} and use~\eqref{mq7ny8m},~\eqref{t0khqm8},~\eqref{nly5agg} and~\eqref{g20r5hg}.
\end{proof}
\begin{theorem}
If $r$ is a complex number that is not a half-integer, then
\begin{align}\label{yoctbpp}
&\sum_{k = 0}^\infty  {\frac{1}{{4^k }}\frac{{\left( {\frac{1}{2}} \right)_k^3 }}{{k!^3 }}\frac{{\binom{{2k + 2r}}{{2r}}}}{{\binom{{k + r}}{r}^2 \binom{{k - 1/2}}{r}}}\left( {\left( {6k + 2r + 1} \right)\left( {H_{k - r - 1/2}  - H_{k + r - 1/2}  + H_{k + r} } \right) - 2} \right)}\nonumber\\ 
&\qquad = \frac{{r!^2 }}{{\left( {\frac{1}{2}} \right)_r^2 }}\,\frac{{H_{ - r - 1/2} }}{{\cos \left( {\pi r} \right)}}\,\frac{4}{\pi }.
\end{align}
\end{theorem}
\begin{proof}
Differentiate~\eqref{k1x81g9} with respect to $c$ to obtain, after some rearrangement,
\begin{align}
&\sum_{k = 0}^\infty  {\frac{{\left( {\frac{1}{2}} \right)_k \left( c \right)_k \left( {1 - c} \right)_k }}{{4^kk!^2 \left( {\frac32 - c} \right)_k }}\left( {\frac{{3k - c + 1}}{2}\left( {H_{c + k - 1}  - H_{ - c + k}  + H_{1/2 - c + k} } \right) - \frac{1}{2}} \right)}\nonumber\\
&\qquad  = \frac{{H_{c - 1} }}{{\sqrt \pi  }}\left( {1 - c} \right)_{1/2} .
\end{align}
Set $c=1/2-r$ to get
\begin{align}\label{pehnhlp}
&\sum_{k = 0}^\infty  {\frac{{\left( {\frac{1}{2}} \right)_k \left( { - r + \frac{1}{2}} \right)_k \left( {r + \frac{1}{2}} \right)_k }}{{4^kk!^2 \left( {r + 1} \right)_k }}\left( {\frac{1}{2}\left( {3k + r + \frac{1}{2}} \right)\left( {H_{k - r - 1/2}  - H_{k + r - 1/2}  + H_{k + r} } \right) - \frac{1}{2}} \right)}\nonumber\\
&\qquad  = \frac{{H_{ - r - 1/2} }}{{\sqrt \pi  }}\left( {r + \frac{1}{2}} \right)_{1/2} ;
\end{align}
and hence~\eqref{yoctbpp} after using~\eqref{l5fmdri} and~\eqref{bjv8z0f}.
\end{proof}
\begin{corollary}
If $n$ is a non-negative integer, then
\begin{align*}
&\sum_{k = 0}^\infty  {\frac{{\left( {\frac{1}{2}} \right)_k^3 }}{{4^k k!^3 }}\frac{{\prod_{j = 1}^{2n} {\left( {2k + j} \right)} }}{{\prod_{j = 1}^n {\left( {k + j} \right)^2 \left( {2k - 2j + 1} \right)} }}\left( {\left( {6k + 1 + 2n} \right)\left( {2O_{k - n}  - 2O_{k + n}  + H_{k + n} } \right) - 2} \right)}\\
&\qquad  =( - 1)^n\,\frac{{2^{n + 3} }}{{\left( {\frac{1}{2}} \right)_n }}\left( {O_n  - \ln 2} \right)\frac{1}{\pi }.
\end{align*}
In particular,
\begin{equation}
\sum_{k = 0}^\infty  {\frac{{\left( {\frac{1}{2}} \right)_k^3 }}{{4^k k!^3 }}\left( {\left( {6k + 1} \right)H_k  - 2} \right)}  =  - \frac{{8\ln 2}}{\pi }.
\end{equation}
\end{corollary}
\begin{proof}
Set $r=n$ in~\eqref{pehnhlp}, use
\begin{equation*}
H_{k - n - 1/2}  - H_{k + n - 1/2}  = 2\left( {O_{k - n}  - O_{k + n} } \right)
\end{equation*}
and identities~\eqref{mq7ny8m} and~\eqref{t0khqm8}.
\end{proof}
\begin{thebibliography}{99}

\bibitem{bailey35} W. N. Bailey, \emph{Generalized Hypergeometric Series}, Cambridge University Press, (1935).

\bibitem{bauer} G. Bauer, Von den Coefficienten der Reihen von Kugelfunktionen einer Variabeln, \emph{Journal f\"ur die reine und angewandte Mathematik} {\bf 56} (1859), 101--121.

\bibitem{campbell24} J. M. Campbell and P. Levrie, The Bauer-Ramanujan formula: historical analyses and perspectives, \emph{British Journal for the History of Mathematics} {\bf 39}:3 (2024), 193--214.

\bibitem{campbell26} J. M. Campbell, Applications of a class of transformations of complex sequences, \emph{Journal of the Ramanujan Mathematical Society}, to appear .

\bibitem{chu14} W. Chu and W. Zhang, Accelerating Dougall's  ${}_5F_4$-sum and infinite seriesinvolving $\pi$, \emph{Mathematics of Computation} {\bf 285} (2014), 475--512.

\bibitem{dougall06} J. Dougall, Vandermonde's theorem and some more general expansions, \emph{Proceedings of the Edinburgh Mathematical Society} {\bf 25} (1906), 114--152.

\bibitem{ekhad94} S. B. Ekhad and D. Zeilberger, A $WZ$ proof of Ramanujan's formula for $\pi$, in J. M. Rassias, ed., \emph{Geometry, Analysis and Mechanics}, World Scientific, 1994, pp. 107--108.

\bibitem{hou24} Q. Hou, H. He and X. Wang, Ramanujan-inspired series for $1/\pi$ involving harmonic numbers, \emph{Journal of Differential Equations of Applications} {\bf 30}:6 (2024), 775--786.

\bibitem{levrie10} P. Levrie, Using Fourier-Legendre expansions to derive series for $1/\pi$ and $1/\pi^2$, \emph{Ramanujan Journal} {\bf 22} (2010), 221--230.

\bibitem{ramanujan14} S. Ramanujan, Modular equations and approximations for $\pi$, \emph{Quarterly Journal of Pure and Applied Mathematics} {\bf 45} (1914), 350--372.

\end{thebibliography}

\end{document}
